\documentclass[11pt,a4paper]{article}
\usepackage[T1]{fontenc}\usepackage[utf8]{inputenc}\usepackage{lmodern}
\usepackage[a4paper,margin=2.2cm]{geometry}
\usepackage{amsmath,amssymb}\usepackage{enumitem}\usepackage{booktabs,longtable,array}
\usepackage[hidelinks]{hyperref}
\setlength{\parindent}{0pt}\setlength{\parskip}{0.4em}
\title{\textbf{OP-CONST T1b/T1c: selective dressing, adiabatic pinning, and a
conditional protection theorem}\\\large Recon v1 --- proposal-only; includes the
citation pass for the stress anchors}
\author{Per reviewer roadmap after D1 (items: T1b $u_0$ anchoring; T1c lemma
promotion; anchor verification)}
\date{12 June 2026}
\begin{document}\maketitle

\section*{0. Scope and firewall}
This document closes the two physics nodes named by the review --- T1b (anchoring
of the drifting amplitude) and T1c (promotion/disposition of the $S_0$ candidate
lemma) --- and performs the citation pass for the varying-constants stress
anchors. Everything is proposal-layer; no preprint edits. The outcome is a
\emph{conditional protection theorem at the recorded-closure level}, with each
hypothesis mapped to an already-registered open problem, plus one honest
correction to the earlier $S_0$-lever emphasis.

\section{T1b --- the recorded coupling architecture decides the anchoring
question}
\textbf{(a) What the record says.} The $\phi$-first IR closure
(eq:phi\_first\_action, L$\sim$18438) is
\[
S_{\rm IR}=\int d^4x\sqrt{-g}\Big\{\tfrac{\bar M_{\rm Pl}^2}{2}e^{\beta_\phi\phi}R
-\tfrac{\bar M_{\rm Pl}^2}{2}\omega(\phi)(\partial\phi)^2-U(\phi)
+\mathcal L_{\rm mat}[g_{\mu\nu},\Psi]\Big\},
\]
with $e^{\beta_\phi\phi}=u/u_\infty$ (eq:phi\_def) and A4 (``slow matter sectors
couple to $g^{\rm eff}_{\mu\nu}$''; Level~B, OP-matter). The cosmological
amplitude excursion $\phi(X)$ dresses \emph{only} the Einstein--Hilbert term.
$\mathcal L_{\rm mat}$ carries no $\phi$ argument: a direct search of the closure
finds no $e^{\#\phi}F_{AB}F^{AB}$ operator, no $\phi$-dependent Yukawas, and no
$\phi$ entering the $v_2$/$\mathcal I_{\rm EW}$ chain.

\textbf{(b) Classification.} This is the Jordan-frame Brans--Dicke-class
architecture (scalar couples to the gravitational sector only), \emph{not} the
Bekenstein class (scalar coupled to $F^2$). The standard consequence applies:
with universal metric coupling and no direct $\phi$--matter operators, all
dimensionless matter couplings are constants of $\mathcal L_{\rm mat}$;
\[
\zeta_A=c_{v_2}=0\qquad\text{identically at the recorded-closure level,}
\]
and the only quantity that varies cosmologically is $G_{\rm eff}(\phi)=G_N
e^{-\beta_\phi\phi}$ --- precisely the single channel for which the preprint
already records a screened/cosmological split and LLR compliance.

\textbf{(c) Resolution of the Steps 1--3 tension.} The B-tracking stress
scenarios of Steps 1--3 (power-law or endpoint-log tracking of $Z_A$, $v_2$)
correspond, in this language, to completion-layer operators of the form
$e^{\gamma_i\beta_\phi\phi}\,\mathcal O_i$ inside $\mathcal L_{\rm mat}$ ---
operators that are \emph{absent from the recorded closure}. The matching
coefficients ($K_\theta^{\rm eff}$, $\kappa_n$, $Z^{(0)}$, the $v_2$ chain) enter
$\mathcal L_{\rm mat}$ as constants constructed at the parameter point; the
cosmological excursion is routed exclusively into the $R$-dressing by the
closure's construction. OP-CONST therefore reduces to: \emph{does the microscopic
completion preserve this selective coupling?} --- which is the OP-grad layer,
now with quantitative thresholds (\S5).

\section{T1c(i) --- adiabatic pinning of the $\Phi$-sector (derived)}
The potential-sector amplitude is gapped at $m_\varphi=\sqrt{2\lambda}\,\phi_0$
with $\phi_0\simeq\bar M_{\rm Pl}$. A heavy field tracking a slowly varying
minimum is displaced by $\delta u_\Phi/\phi_0\sim(H/m_\varphi)^2$. Numerically,
\[
\Big(\frac{H_0}{m_\varphi}\Big)^2=\frac{(1.44\times10^{-42}\,{\rm GeV})^2}
{2\lambda\,(2.435\times10^{18}\,{\rm GeV})^2}\;\approx\;\frac{1.7\times10^{-121}}
{\lambda},
\]
i.e.\ $\lesssim10^{-120}$ for any $\lambda\gtrsim10^{-2}$, and still
$\sim10^{-110}$ at matter--radiation equality. \textbf{Consequence:} the
$\Phi$-sector evaluation point of the matching formulas
($K_\theta^{\rm eff}\sim\phi_0^2$, $m_\varphi^2=2\lambda\phi_0^2$,
$v_2=\phi_0e^{-\mathcal I_{\rm EW}}$) is parameter-anchored to a precision some
$115$ orders of magnitude beyond the required thresholds. The
``same-evaluation-context'' assumption of the T1 candidate lemma closes at the
recorded level, and the $\alpha$/$\mu$ chains cannot be fed by $\Phi$-sector
excursions.

\section{T1c(ii) --- honest correction: the $S_0$-lever protects rescaling, not
shape excursions}
Explicit check in the simplest quartic, $V=\tfrac{\lambda}{4}(u^2-\phi_0^2)^2$,
with $K_\theta=u^2$ and radial gap $m^2(u)=\lambda(3u^2-\phi_0^2)$:
\[
\frac{d\ln(K_\theta/m^2)}{d\ln u}\Big|_{u=\phi_0}
=2-\frac{6u^2}{3u^2-\phi_0^2}\Big|_{u=\phi_0}=2-3=-1 .
\]
The ratio $K_\theta/m_\varphi^2$ is therefore \emph{not} invariant under
off-minimum excursions: the $S_0$-lever cannot by itself be promoted to the
protection theorem. Its corrected role is twofold: (i)~consistency of the
recorded matching at the pinned point ($S_0^{\rm EFT}=c_\perp/4\lambda$,
parameter-only); (ii)~robustness against \emph{common rescalings} of the
$\Phi$-sector normalisation (the T1 candidate no-go lemma survives in this
rescaling form). The primary shield is the pinning of \S2 together with the
selective coupling of \S1 --- not the ratio identity.

\section{Conditional protection theorem (recorded-closure level)}
\textbf{Hypotheses.}
\begin{description}[leftmargin=1.9em,itemsep=2pt]
\item[H1.] (Recorded closure.) The cosmological scalar dresses the gravitational
sector only, as written in eq:phi\_first\_action. \emph{Status: established by
the record as written.}
\item[H2.] (Metric universality.) Slow matter couples universally to
$g^{\rm eff}_{\mu\nu}$ (A4). \emph{Status: Level~B; this is OP-matter.}
\item[H3.] (Selective-coupling completion.) The microscopic completion generates
no $e^{\gamma_i\beta_\phi\phi}\mathcal O_i$ operators in $\mathcal L_{\rm mat}$
with $|\gamma_A|\gtrsim3\times10^{-5}$ (EM kinetic sector) or
$|\gamma_\mu|\gtrsim3\times10^{-6}$ ($v_2$/mass chain). \emph{Status: open; this
is the OP-grad layer, now quantified.}
\item[H4.] (Adiabatic pinning.) $\Phi$-sector displacement
$\lesssim(H/m_\varphi)^2\sim10^{-120}$. \emph{Status: derived (\S2).}
\end{description}
\textbf{Statement.} Under H1--H4, on the cosmological branch of the retained
$\varepsilon$-band: $\Delta\alpha/\alpha=\Delta\mu/\mu=0$ up to the H3-bounded
completion couplings, for all redshifts; the quasar/H$_2$/clock channels are
satisfied identically; and the only varying coupling is $G_{\rm eff}$, with the
already-recorded screened/cosmological split. \textbf{Proof sketch.} H1+H2 put
all dimensionless couplings into the $\phi$-independent $\mathcal L_{\rm mat}$
(Brans--Dicke-class statement); H4 anchors the matching point of the
$\mathcal L_{\rm mat}$ coefficients; H3 bounds the only remaining channel
(completion operators). $\square$

\textbf{Corollary (architecture falsifier).} A confirmed cosmological detection
of $\Delta\alpha/\alpha$ at the ppm level within the retained band would imply
$|\gamma_A|\sim2\times10^{-5}$, i.e.\ falsify selective coupling (H3) or the band
itself --- a sharp, two-sided structural test.

\textbf{What the theorem is and is not.} It is a theorem \emph{given the recorded
closure} (H1) --- the protection is not an extra mechanism but a property of the
coupling architecture the preprint already wrote down. It is conditional exactly
where the theory is open: H2 = OP-matter, H3 = OP-grad completion. It does
\emph{not} claim that the microscopic theory has been shown to satisfy H3; it
converts OP-CONST from ``unknown exposure'' into ``two named hypotheses with
quantified thresholds''.

\section{Citation pass --- verified stress anchors (this round)}
\renewcommand{\arraystretch}{1.2}\small
\begin{longtable}{@{}p{3.4cm}p{4.7cm}p{2.2cm}p{3.0cm}@{}}
\toprule \textbf{Channel} & \textbf{Verified bound} & $|\Delta\ln u|$
($\varepsilon_*$; band) & \textbf{Threshold on $\gamma$} \\ \midrule \endhead
\bottomrule \endfoot
Quasar $\alpha$, ESPRESSO single system $z=1.15$ & $\Delta\alpha/\alpha=
1.3\pm1.3_{\rm stat}\pm0.4_{\rm sys}$ ppm; combined $|\Delta\alpha/\alpha|<2.7$
ppm ($3\sigma$) [Murphy et al.\ 2022, A\&A 658, A123] & $0.049$
($0.045$--$0.057$) & $|\gamma_A|\lesssim5\times10^{-5}$ ($3\sigma$) \\
Quasar $\alpha$, weighted mean $z=0.6$--$2.4$ (29 mitigated meas.) &
$\langle\Delta\alpha/\alpha\rangle=-0.5\pm0.5_{\rm stat}\pm0.4_{\rm sys}$ ppm
[ibid.] & $0.05$--$0.07$ & $|\gamma_A|\lesssim(2$--$3)\times10^{-5}$
($2\sigma$) \\
Radio $\mu$ (CH$_3$OH), $z=0.886$ & $|\Delta\mu/\mu|\le1.1\times10^{-7}$
($2\sigma$, stat); robust $\lesssim4\times10^{-7}$ [Kanekar et al.\ 2015, MNRAS
448, L104] & $0.041$ ($0.038$--$0.048$) & $|\gamma_\mu|\lesssim3\times10^{-6}$
($2\sigma$, stat) \\
Radio $\mu$ (NH$_3$), $z=0.685$ & $|\Delta\mu/\mu|<2.4\times10^{-7}$ ($2\sigma$)
[Kanekar 2011] & $0.034$ & $|\gamma_\mu|\lesssim7\times10^{-6}$ \\
H$_2$ $\mu$, $z=2$--$4.2$ (12 systems) & $|\Delta\mu/\mu|<10^{-5}$; best systems
few$\times10^{-6}$ [e.g.\ Dapr\`a et al.\ 2017, MNRAS 465, 4057] & $0.071$--
$0.106$ & $|\gamma_\mu|\lesssim6\times10^{-5}$ \\
Clocks (local), Yb$^+$ E3/E2 & $\dot\alpha/\alpha=1.0(1.1)\times10^{-18}$/yr;
$\dot\mu/\mu=-8(36)\times10^{-18}$/yr [Lange et al.\ 2021, PRL 126, 011102] &
screened + undressed & doubly protected; consistency channel \\
\end{longtable}\normalsize
Corrections relative to the Step-1--3 working anchors: the $10^{-7}$-level $\mu$
constraint sits at $z<1$ (radio molecules), not at $z\simeq2.5$; the H$_2$
channel at $z\sim2.4$ is $\sim5\times10^{-6}$. The binding completion thresholds
become $|\gamma_A|\lesssim(2$--$3)\times10^{-5}$ and
$|\gamma_\mu|\lesssim3\times10^{-6}$ --- same orders as the recon estimates, now
anchored to named references. Oklo and meteorite bounds are $1$--$2$ orders
weaker than clocks and are likewise local/screened channels.

\section{Status and residue}
\textbf{Closed this round:} T1b (the anchoring question is decided \emph{at the
recorded-closure level} by the selective-dressing architecture: $\zeta_A=
c_{v_2}=0$ identically); T1c (pinning derived at $10^{-120}$; $S_0$-lemma
honestly repositioned to its rescaling-robustness role); citation pass (all
anchors verified with named references; thresholds updated). \textbf{Open
residue for the unconditional statement:} H2 (OP-matter first-principles
universality) and H3 at the microscopic level --- deriving that the OP-grad
completion generates no above-threshold $e^{\gamma\phi}$ matter operators. This
residue \emph{is} the completion programme itself; no varying-constants-specific
unknown remains. \textbf{Proposed next D1 (pending review):} upgrade the
preprint OP-CONST block from ``open requirement'' to ``conditional protection
route'': recorded-closure theorem statement (H1--H4), the falsifier corollary,
verified anchors with bibliography entries, and the corrected disposition of the
$S_0$ identity. No preprint contact before that review.
\end{document}
